% Generated draft: requires mathematical and source review.
\documentclass[11pt]{article}
\usepackage[margin=25mm]{geometry}
\usepackage{fontspec}
\setmainfont{Latin Modern Roman}
\usepackage{amsmath,amssymb,mathtools,unicode-math}
\setmathfont{Latin Modern Math}
\usepackage{enumitem,needspace,xurl,hyperref}
\setlength{\parindent}{0pt}
\setlength{\parskip}{5pt}
\setlength{\emergencystretch}{3em}
\newcommand{\N}{\mathbb{N}}
\newcommand{\Z}{\mathbb{Z}}
\newcommand{\Q}{\mathbb{Q}}
\newcommand{\R}{\mathbb{R}}
\newcommand{\C}{\mathbb{C}}
\newcommand{\F}{\mathbb{F}}
\newcommand{\A}{\mathbb{A}}
\newcommand{\PP}{\mathbb{P}}
\newcommand{\E}{\mathbb{E}}
\newcommand{\eps}{\varepsilon}
\newcommand{\dd}{\,\mathrm{d}}
\begin{document}
\section*{UnsolvedMath ID: 30003624}
\textbf{Source problem code:} catalogue entry
\subsection{Definitions and mathematical statement}
The supplied record concerns a ``Toda shock problem.'' Write \(n\) and \(t\) for its spatial and temporal variables, respectively, and retain \(\xi_{\rm cr}\) and \(\xi_{\rm cr,1}\) as the source's symbols for the critical ray parameters. The governing Toda equations, the domains of these variables, the initial and boundary conditions, and the definitions and values of the critical parameters are not supplied. The ratios below require \(t\ne 0\).

The source asks: determine the long-time asymptotics for the Toda shock problem in transitional regions near the two rays
\[
\frac{n}{t}=\xi_{\rm cr},\qquad
\frac{n}{t}=\xi_{\rm cr,1},
\]
including any asymptotic solitons that may occur.

This is an incompletely specified asymptotic problem. The record does not identify the quantities whose asymptotics are sought, the direction of the long-time limit, the scaling or widths of the transitional regions, the admissible class of shock data, or the meaning and conditions for occurrence of asymptotic solitons. Consequently, it does not determine a precise quantified asymptotic assertion, an error estimate, or a uniformity regime. These specifications must be recovered from an authoritative formulation before the question can be made mathematically self-contained.
\subsection{Short English statement}
Determine the long-time behavior of the unspecified Toda shock problem near the critical rays \(n/t=\xi_{\rm cr}\) and \(n/t=\xi_{\rm cr,1}\), including possible asymptotic solitons. The supplied record lacks the equations, data, critical parameters, observables, and transition scaling needed for a precise formulation.
\subsection{Sources}
\begin{itemize}
\item Supplied local UnsolvedMath registry record, permanent ID \(30003624\), title ``Transitional Asymptotics for Toda Shock Waves,'' source problem code \texttt{catalogue entry}, in the Oberwolfach Reports Open Problems collection. The recorded proposed year is 2017; no individual proposer is identified.
\item Supplied source citation: \emph{Mini-Workshop: Reflectionless Operators: The Deift and Simon Conjectures (2018)}, Oberwolfach Reports, pp.\ 2943--2985, DOI \texttt{10.4171/owr/2017/49}. The literature-source label is ``Oberwolfach Report 49/2017.'' Supplied link: \url{https://doi.org/10.4171/owr/2017/49}. The record describes these reports as documenting mathematical workshops and their open-problem discussions. The relationship between the citation's 2018 date and its 2017 report designation is not explained.
\item Literature assessment dated exactly \texttt{2026-08-21}: the supplied curation labels the problem open, while stating that the extracted record names transitional asymptotics without specifying the quantity or regime, so no definitive resolution can be verified. It lists Oberwolfach Report 49/2017 as the source checked. This is supplied dated curation, not an independently verified literature assessment; its status and cited claims require independent checking.
\item The supplied statement status is \texttt{corrected\_verified}, accompanied by a claim that mathematical content was preserved while notation, typography, or omitted definitions were repaired. This is a registry annotation, not independent verification of the formulation or its completeness.
\end{itemize}
\end{document}
